% Generated draft: requires mathematical and source review.
\documentclass[11pt]{article}
\usepackage[margin=25mm]{geometry}
\usepackage{fontspec}
\setmainfont{Latin Modern Roman}
\usepackage{amsmath,amssymb,mathtools,unicode-math}
\setmathfont{Latin Modern Math}
\usepackage{enumitem,needspace,xurl,hyperref}
\setlength{\parindent}{0pt}
\setlength{\parskip}{5pt}
\setlength{\emergencystretch}{3em}
\newcommand{\N}{\mathbb{N}}
\newcommand{\Z}{\mathbb{Z}}
\newcommand{\Q}{\mathbb{Q}}
\newcommand{\R}{\mathbb{R}}
\newcommand{\C}{\mathbb{C}}
\newcommand{\F}{\mathbb{F}}
\newcommand{\A}{\mathbb{A}}
\newcommand{\PP}{\mathbb{P}}
\newcommand{\E}{\mathbb{E}}
\newcommand{\eps}{\varepsilon}
\newcommand{\dd}{\,\mathrm{d}}
\begin{document}
\section*{UnsolvedMath ID: 701777}
\textbf{Source problem code:} catalogue entry
\subsection{Definitions and mathematical statement}
Let \(p\) be a prime. A pro-\(p\) group is a topological group isomorphic to an inverse limit of finite \(p\)-groups. It is finitely presented as a pro-\(p\) group if it is isomorphic to \(F/R\), where \(F\) is a free pro-\(p\) group on finitely many generators and \(R\) is the closed normal subgroup generated by finitely many elements.

A finite connected graph of pro-\(p\) groups \(\mathcal G\) consists of a finite connected graph \(\Gamma\), pro-\(p\) vertex and edge groups, and continuous injective homomorphisms from each edge group to its incident vertex groups. Write \(\Pi_1^{(p)}(\mathcal G)\) for its fundamental group in the category of pro-\(p\) groups. The graph of groups is proper if the canonical maps from its vertex groups to \(\Pi_1^{(p)}(\mathcal G)\) are injective. We use the standard reducedness convention that, for every edge joining distinct vertices, neither of its two edge-to-vertex homomorphisms is surjective; loop edges are exempt from this condition. Edges are counted without orientation. The supplied record does not explicitly define these graph-of-groups conventions, so their agreement with the cited source requires review.

A pro-\(p\) group \(G\) is called accessible if there exists \(n(G)\in\Z_{\geq 0}\) such that every finite proper reduced graph of pro-\(p\) groups \(\mathcal G\) with finite edge groups and with \(\Pi_1^{(p)}(\mathcal G)\cong G\) as topological groups satisfies
\[
\lvert E(\Gamma)\rvert\leq n(G).
\]
Here the finite edge groups are necessarily finite \(p\)-groups, and the bound may depend on \(G\), but must be independent of its graph-of-groups decomposition.

\textbf{Question.} For every prime \(p\), is every finitely presented pro-\(p\) group accessible?

The supplied registry labels this question open; that label is not an independently verified mathematical or literature assessment.
\subsection{Short English statement}
Does every finitely presented pro-\(p\) group admit a bound, depending only on the group, on the number of edges in all its finite proper reduced graph-of-pro-\(p\)-groups decompositions with finite edge groups? The graph-of-groups conventions and the attribution discrepancy in the supplied record require review.
\subsection{Sources}
\begin{itemize}
\item Supplied local UnsolvedMath registry record, permanent ID \(701777\), exact source problem code \texttt{catalogue entry}, titled \emph{Kourovka Notebook Problem 20.119}. The record supplies status \texttt{open} and was updated on \texttt{2026-10-16T14:00:00Z}; this is supplied registry information, not independent verification.
\item \emph{The Kourovka Notebook (Unsolved Problems in Group Theory)}, problem 20.119, attributed in the supplied background to the 20th issue (2022) and the 21st issue edition (2026), arXiv:1401.0300v48. The supplied collection description names E. I. Khukhro and V. D. Mazurov as editors. These bibliographic details have not been independently checked.
\item The statement cites \emph{J. Algebra} \textbf{525} (2019), 1--18, in connection with the accessibility definition. No article title, author list, DOI, or URL is supplied.
\item Attribution as supplied: the statement begins ``(G. Wilkes)'', while the background and proposer field name P. Zalesskii, with proposed year 2022. The record does not explain the relationship between these attributions.
\end{itemize}
\end{document}
